% ============================================================
% sec/04metodologia.tex
%
% OBLIGATORIO: formalización con formulación matemática,
% pipeline previsto, protocolo experimental con particiones,
% validación y métricas justificadas, plan de baselines, y
% riesgos con su mitigación.
%
% Cada decisión sin justificación técnica se penaliza.
% ============================================================
\section{Metodología Propuesta}
\label{sec:metodologia}

\subsection{Formalización del problema}
\label{subsec:formalizacion}

El problema se formula como clasificación binaria probabilística sobre
datos tabulares. Sea $\mathcal{E}$ el conjunto de episodios disponibles,
donde cada episodio $e \in \mathcal{E}$ corresponde a una partida completa
entre dos agentes. Cada episodio produce una secuencia de estados
observables $\{s_{e,1}, s_{e,2}, \dots, s_{e,T_e}\}$, donde $T_e$ denota la
cantidad de estados registrados en el episodio $e$, y un desenlace
$y_e \in \{0,1\}$ que indica si el jugador de referencia resultó vencedor.

Cada estado se transforma en un vector de características
$\mathbf{x}_{e,t} = \phi(s_{e,t}) \in \mathbb{R}^d$ mediante una función de
extracción $\phi$ restringida a la información disponible en el instante
$t$. La etiqueta asociada a toda observación del episodio $e$ es el mismo
valor $y_e$.

El objetivo es estimar
\begin{equation}
\label{eq:objetivo}
\hat{p}(\mathbf{x}) \approx \Pr\!\left(Y = 1 \mid \mathbf{X} = \mathbf{x}\right),
\end{equation}
donde la probabilidad se toma sobre la aleatoriedad del juego posterior al
instante observado: el orden restante del mazo, la información oculta del
oponente y las decisiones futuras de ambos agentes.

\razon{La ecuación debe explicitar sobre qué se toma la probabilidad. Es lo
que distingue este problema de uno determinista, y es la respuesta al
criterio de exclusión que Grinsztajn \textit{et al.} aplican a los conjuntos
de juegos en los cuales el objetivo se deduce de los datos
\cite{grinsztajn_trees_2022}.}

\subsubsection{Unidad de observación y unidad independiente}

Esta distinción gobierna la totalidad del diseño experimental y se declara
de forma explícita:

\begin{itemize}
  \item \textbf{Unidad de observación:} el estado $\mathbf{x}_{e,t}$. Es la
        fila de la tabla y la entrada del modelo.
  \item \textbf{Unidad estadísticamente independiente:} el episodio $e$. El
        desenlace se sortea una única vez por episodio, y todas las
        observaciones de un mismo episodio comparten ese sorteo.
\end{itemize}

En consecuencia, el número de observaciones excede en más de dos órdenes de
magnitud el número de sorteos independientes del desenlace. Brill
\textit{et al.} demuestran que, bajo esta estructura, el tamaño de muestra
efectivo resulta una fracción del nominal, y que la precisión del estimador
decrece conforme aumenta el número de observaciones por unidad manteniendo
constante el total de filas \cite{brill_winprob_2026}.

\razon{Toda mención del tamaño del conjunto de datos en este documento debe
expresar ambas cifras: el número de observaciones y el número de episodios.
Reportar únicamente el conteo de filas transmitiría una noción inexacta de
la cantidad de evidencia disponible, precisamente el error que Brill
\textit{et al.} identifican \cite{brill_winprob_2026}.}

\subsubsection{Definición del objetivo}

\pc{Descripción de cómo se deriva la etiqueta a partir de la información
del episodio, y qué episodios se excluyen por no admitir una derivación
inequívoca. Corresponde a la auditoría del conjunto de datos, Eje C.}

\subsection{Alcance del uso de las observaciones}
\label{subsec:alcanceobs}

Se conserva la totalidad de las observaciones por episodio, sin
submuestreo. Brill \textit{et al.} reportan que retener todas las
observaciones de una unidad resulta preferible a conservar una sola, dado
que el conjunto completo aporta información sobre la estructura del espacio
de covariables, pese a la dependencia entre las etiquetas
\cite{brill_winprob_2026}.

\razon{Esta decisión debe justificarse de forma explícita porque la
alternativa ---retener un estado por episodio para obtener observaciones
independientes--- resulta intuitivamente atractiva y es peor. Sin la cita,
la decisión parecería una omisión en lugar de una elección informada.}

\subsection{Pipeline previsto}
\label{subsec:pipeline}

El procedimiento se organiza en seis etapas:

\begin{enumerate}
  \item \textbf{Ingesta.} Lectura de las repeticiones en su formato
        original.
  \item \textbf{Extracción de estados.} Recorrido de cada episodio y
        obtención de los estados observables, junto con el desenlace.
  \item \textbf{Construcción de características.} Aplicación de $\phi$
        sobre cada estado, restringida a la información disponible en el
        instante correspondiente.
  \item \textbf{Materialización.} Escritura de la tabla resultante en
        formato columnar, ejecutada una sola vez.
  \item \textbf{Partición.} División agrupada por episodio en los
        subconjuntos descritos en la Sección~\ref{subsec:protocolo}.
  \item \textbf{Entrenamiento y evaluación.} Ajuste de los modelos de
        referencia y medición bajo los tres protocolos de generalización.
\end{enumerate}

Xenopoulos \textit{et al.} describen un procedimiento análogo, en el cual
las repeticiones se procesan mediante un analizador sintáctico que produce
estados del juego a intervalos regulares, y dichos estados constituyen la
entrada de los modelos \cite{xenopoulos_winprob_2022}. Hodge \textit{et al.}
emplean igualmente un analizador sobre repeticiones para producir vectores
de estado por instante \cite{hodge_winprediction_2021}.

\razon{Citar ambos precedentes evita que el procedimiento parezca una
construcción improvisada. La estructura repetición--analizador--tabla es la
práctica establecida del dominio.}

% ============================================================
% Figura: pipeline propuesto.
% Dibujada con TikZ: no requiere archivo externo.
% Las seis etapas corresponden a la enumeración de esta sección.
% ============================================================
\begin{figure}[!t]
\centering
\begin{tikzpicture}[
  font=\scriptsize,
  etapa/.style={
    draw=tecrule, fill=boxbg, thick, rounded corners=2pt,
    text width=5.4cm, align=center, inner sep=4pt, minimum height=7mm},
  hito/.style={
    draw=tecnavy, fill=tecnavy!8, thick, rounded corners=2pt,
    text width=5.4cm, align=center, inner sep=4pt, minimum height=7mm},
  flecha/.style={-{Stealth[length=4pt]}, draw=tecsteel, thick},
  nota/.style={font=\tiny\itshape, text=labelgray, align=left,
               text width=1.9cm},
  fase/.style={font=\tiny\bfseries, text=tecnavy, rotate=90,
               align=center}
]

\node[etapa] (e1) {1. Ingesta de repeticiones};
\node[etapa, below=3.5mm of e1] (e2) {2. Extracción de estados y desenlace};
\node[etapa, below=3.5mm of e2] (e3) {3. Construcción de características $\phi$};
\node[hito,  below=3.5mm of e3] (e4) {4. Materialización de la tabla};
\node[hito,  below=3.5mm of e4] (e5) {5. Partición agrupada por episodio};
\node[etapa, below=3.5mm of e5] (e6) {6. Entrenamiento y evaluación};

\draw[flecha] (e1) -- (e2);
\draw[flecha] (e2) -- (e3);
\draw[flecha] (e3) -- (e4);
\draw[flecha] (e4) -- (e5);
\draw[flecha] (e5) -- (e6);

\node[nota, right=2.5mm of e3] {solo información disponible en el instante $t$};
\node[nota, right=2.5mm of e4] {ejecutada una sola vez};
\node[nota, right=2.5mm of e5] {ningún episodio en dos subconjuntos};
\node[nota, right=2.5mm of e6] {protocolos A, B y C};

\begin{scope}[on background layer]
  \node[draw=tecrule, dashed, rounded corners=3pt, inner sep=3pt,
        fit=(e1)(e4)] (faseA) {};
  \node[draw=tecrule, dashed, rounded corners=3pt, inner sep=3pt,
        fit=(e5)(e6)] (faseB) {};
\end{scope}

\node[fase, left=1mm of faseA] {Construcción\\del conjunto};
\node[fase, left=1mm of faseB] {Experimentación};

\end{tikzpicture}
\caption{Pipeline propuesto. Las etapas 1 a 4 se ejecutan una sola vez y
producen la representación tabular; las etapas 5 y 6 operan sobre ella. La
partición de la etapa 5 asigna episodios completos a un único subconjunto,
condición que se verifica de forma automática.}
\label{fig:pipeline}
\end{figure}

\subsection{Protocolo experimental}
\label{subsec:protocolo}

\subsubsection{Partición agrupada por episodio}

Toda partición se realiza asignando episodios completos a un único
subconjunto. Ninguna observación de un episodio puede figurar en dos
subconjuntos distintos.

Esta decisión responde a dos evidencias convergentes. Brill \textit{et al.}
muestran que la dependencia entre observaciones que comparten desenlace
incrementa el sesgo y la varianza del estimador
\cite{brill_winprob_2026}. Figueiredo y Mendes cuantifican el efecto sobre
las métricas reportadas: la partición aleatoria sobre observaciones
altamente correlacionadas infla el desempeño medido respecto de una
partición que asigna agrupaciones completas
\cite{figueiredo_leakage_2024}.

\razon{Se citan las dos referencias y no una, porque miden efectos
distintos: Brill sobre las propiedades del estimador, Figueiredo sobre la
cifra que se publica. Presentar solo una dejaría el argumento a medias.}

Una ventaja de este conjunto de datos respecto de los escenarios estudiados
por Figueiredo y Mendes es que la agrupación no requiere inferencia: el
identificador de episodio la determina de forma exacta. En su trabajo, la
contribución central consiste precisamente en un procedimiento para inferir
agrupaciones que no vienen dadas \cite{figueiredo_leakage_2024}.

\subsubsection{Tres protocolos de generalización}

\begin{table}[!t]
\centering
\caption{Protocolos de evaluación y pregunta que responde cada uno.}
\label{tab:protocolos}
\footnotesize
\begin{tabularx}{\columnwidth}{@{}P{1.5cm} P{2.3cm} Y@{}}
\toprule
\textbf{Protocolo} & \textbf{Criterio de partición} &
\textbf{Pregunta que responde}\\
\midrule
A & Episodios disjuntos, asignación aleatoria &
¿Generaliza a partidas no observadas?\\
B & Episodios disjuntos, corte temporal &
¿Se conserva el desempeño sobre periodos posteriores?\\
C & Episodios disjuntos, mazos disjuntos &
¿Aprendió a leer el estado o memorizó qué configuraciones vencen?\\
\bottomrule
\end{tabularx}
\end{table}

El protocolo B responde a la evidencia de que la población de estrategias
de una comunidad competitiva evoluciona. Saravanan y Guzdial caracterizan
dicho conjunto de estrategias dominantes y documentan que cambia de forma
continua \cite{saravanan_metadiscovery_2024}. Hodge \textit{et al.}
identifican como limitación del trabajo previo que los datos se recolecten
a lo largo de periodos que atraviesan cambios significativos del juego, y
adoptan una partición ordenada cronológicamente con el fin de no emplear
datos posteriores para predecir eventos anteriores
\cite{hodge_winprediction_2021}.

\razon{El protocolo temporal debe justificarse con literatura del dominio de
juegos y no únicamente con referencias metodológicas de otros campos. Hodge
\textit{et al.} y Saravanan y Guzdial cumplen esa función.}

\pc{El protocolo C requiere el análisis de concentración de configuraciones
de mazo, correspondiente al Eje C, para establecer su viabilidad y el
criterio de separación.}

\subsubsection{Conjunto de calibración}

El procedimiento de recalibración descrito en la
Sección~\ref{subsec:planbaselines} requiere un subconjunto de datos no
empleado en el ajuste del modelo. Silva Filho \textit{et al.} establecen
que la recalibración debe ajustarse sobre datos no vistos durante el
entrenamiento \cite{silvafilho_calibration_2023}. Niculescu-Mizil y Caruana
detallan la consecuencia de incumplirlo: si el modelo separa perfectamente
el conjunto sobre el cual se calibra, la transformación aprendida degenera
\cite{niculescumizil_probabilities_2005}.

\decidir{El diseño requiere resolver de qué subconjunto procederán los
datos de calibración. Las dos alternativas son: una cuarta partición
independiente, obtenida dividiendo el conjunto de entrenamiento y agrupando
por episodio; o una validación cruzada agrupada por identificador de
episodio. La primera resulta más sencilla de auditar; la segunda aprovecha
la totalidad de los datos, pero exige verificar que el procedimiento de
validación cruzada respete la agrupación, dado que las implementaciones
habituales asumen independencia entre observaciones. La decisión debe
tomarse antes de ejecutar cualquier entrenamiento, puesto que modificar las
particiones a mitad del proceso invalidaría los resultados previos.}

\subsubsection{Métricas y su justificación}

\begin{table}[!t]
\centering
\caption{Métricas de evaluación y componente que mide cada una.}
\label{tab:metricas}
\footnotesize
\begin{tabularx}{\columnwidth}{@{}P{2.1cm} P{1.9cm} Y@{}}
\toprule
\textbf{Métrica} & \textbf{Componente} &
\textbf{Justificación}\\
\midrule
Área bajo la curva ROC & Refinamiento &
Mide exclusivamente el ordenamiento. Es invariante ante transformaciones
monótonas de las probabilidades\\
Error de calibración & Calibración &
Mide la correspondencia entre probabilidades predichas y frecuencias
observadas\\
Pérdida logarítmica & Ambos &
Regla de puntuación propia. Penaliza severamente la confianza equivocada\\
Puntaje de Brier & Ambos &
Regla de puntuación propia acotada, con penalización cuadrática\\
\bottomrule
\end{tabularx}
\end{table}

La adopción simultánea de estas cuatro métricas se fundamenta en la
descomposición que Silva Filho \textit{et al.} exponen: las reglas de
puntuación propias se separan en un componente de calibración y uno de
refinamiento \cite{silvafilho_calibration_2023}. Ninguna métrica
individual distingue entre las dos formas en que un estimador puede
fallar, de modo que reportar únicamente una impediría determinar si un
desempeño deficiente obedece a incapacidad de separar las clases o a
probabilidades que no reflejan las frecuencias reales.

La exactitud se descarta como métrica principal. Niculescu-Mizil y Caruana
señalan que, cuando se requieren probabilidades bien calibradas, ni la
exactitud ni el área bajo la curva resultan suficientes
\cite{niculescumizil_probabilities_2005}.

\razon{La justificación de las métricas es uno de los puntos que el
enunciado señala como decisión clave. Debe apoyarse en la descomposición
formal y no en la enumeración de métricas habituales.}

\noindent\textbf{Construcción de las curvas de calibración.} Silva Filho
\textit{et al.} advierten que los estimadores del error de calibración
basados en discretización dependen del número y la construcción de los
intervalos empleados \cite{silvafilho_calibration_2023}. Dimitriadis
\textit{et al.} proponen un procedimiento que produce diagramas
reproducibles junto con una medida numérica de descalibración
\cite{dimitriadis_reliability_2021}.

\decidir{Queda por verificar la disponibilidad de una implementación del
procedimiento de Dimitriadis \textit{et al.} en el entorno de trabajo. En
caso de no estar disponible, se emplea el error de calibración esperado con
discretización, declarando la limitación correspondiente.}

\subsubsection{Intervalos de confianza}

Los intervalos de confianza se construyen mediante remuestreo sobre
episodios y no sobre observaciones individuales. Brill \textit{et al.}
reportan que el remuestreo que ignora la estructura de agrupamiento alcanza
una cobertura marcadamente inferior a la nominal, y que incluso el
remuestreo que respeta dicha estructura permanece por debajo del nivel
declarado \cite{brill_winprob_2026}.

\razon{Debe declararse que los intervalos constituyen una cota optimista de
la incertidumbre real. Presentarlos como intervalos de cobertura nominal
válida contradiría la evidencia de la referencia que fundamenta el método.}

\subsection{Plan de modelos de referencia}
\label{subsec:planbaselines}

\begin{table}[!t]
\centering
\caption{Escalera de modelos de referencia.}
\label{tab:baselines}
\footnotesize
\begin{tabularx}{\columnwidth}{@{}c P{2.4cm} Y@{}}
\toprule
\textbf{Nivel} & \textbf{Modelo} & \textbf{Pregunta que responde}\\
\midrule
B0 & Predictor de clase mayoritaria & ¿Cuál es el piso absoluto?\\
B1 & Heurística de dominio & ¿Se requiere aprendizaje automático?\\
B2 & Regresión logística & ¿La relación es lineal?\\
B3 & Bosques aleatorios & ¿Existen interacciones no lineales?\\
B4 & Potenciación del gradiente & ¿Cuánto aporta la complejidad adicional?\\
\bottomrule
\end{tabularx}
\end{table}

La inclusión de los niveles B2 a B4 se fundamenta en la evidencia de que
los conjuntos de árboles constituyen el estado del arte sobre datos
tabulares \cite{borisov_tabular_2024, grinsztajn_trees_2022}.

El nivel B1 merece atención particular. Hodge \textit{et al.} reportan que
un modelo construido sobre una única característica alcanza un desempeño
cercano al obtenido con el conjunto completo de características
\cite{hodge_winprediction_2021}. Un resultado análogo en este trabajo
constituiría un hallazgo y no un fracaso.

\razon{Declarar de antemano que un buen desempeño de la heurística sería un
resultado válido evita que, ante ese escenario, se busque justificar la
complejidad del modelo con argumentos posteriores a los datos.}

\noindent\textbf{Recalibración.} El procedimiento no aplica recalibración de
forma uniforme. Niculescu-Mizil y Caruana reportan que los métodos de
máximo margen presentan distorsión sistemática, mientras que la regresión
logística y los conjuntos por agregación resultan bien calibrados sin
corrección, y que recalibrar estos últimos puede deteriorarlos
\cite{niculescumizil_probabilities_2005}. En consecuencia, el
procedimiento mide la calibración de cada modelo antes de decidir si
aplicar corrección, y reporta ambas mediciones.

\razon{Los mismos autores reportan que el ordenamiento de los modelos según
la calidad de sus probabilidades difiere antes y después de la
recalibración \cite{niculescumizil_probabilities_2005}. Reportar una sola
de las dos mediciones ofrecería una comparación incompleta.}

\decidir{La elección del método de recalibración permanece abierta.
Niculescu-Mizil y Caruana establecen un criterio basado en el tamaño del
conjunto de calibración \cite{niculescumizil_probabilities_2005}, mientras
que Ojeda \textit{et al.} evalúan sistemáticamente los métodos disponibles
atendiendo, entre otros aspectos, a la generalización hacia poblaciones
externas \cite{ojeda_calibrating_2023}. Dado que el protocolo B de este
trabajo constituye un escenario de transferencia a una población distinta,
corresponde determinar cuál de las dos recomendaciones aplica.}

\subsection{Contrato de reproducibilidad}
\label{subsec:reproducibilidad}

\pc{Semillas fijas y su ubicación, configuraciones almacenadas, registro de
ejecuciones, entorno de cómputo, versiones de dependencias, y el comando
exacto que reproduce cada valor reportado.}

Se adopta además la práctica de repetir cada experimento sobre múltiples
semillas y reportar la dispersión resultante. Grinsztajn \textit{et al.}
repiten su procedimiento de búsqueda variando el orden de las evaluaciones
para separar el efecto del método de la variabilidad del procedimiento
\cite{grinsztajn_trees_2022}. Figueiredo y Mendes ejecutan múltiples ciclos
de entrenamiento por cada estrategia de partición
\cite{figueiredo_leakage_2024}. Brill \textit{et al.} generan múltiples
conjuntos de entrenamiento por configuración \cite{brill_winprob_2026}.

\razon{Los tres precedentes se citan juntos porque establecen la repetición
como práctica del campo y no como precaución particular de este trabajo.}

\subsection{Riesgos y mitigación}
\label{subsec:riesgos}

\begin{table}[!t]
\centering
\caption{Riesgos identificados y su mitigación.}
\label{tab:riesgos}
\footnotesize
\begin{tabularx}{\columnwidth}{@{}P{2.0cm} c Y@{}}
\toprule
\textbf{Riesgo} & \textbf{Impacto} & \textbf{Mitigación}\\
\midrule
Fuga por partición aleatoria & Alto &
Partición agrupada por episodio con verificación automática\\[2pt]
Fuga por preprocesamiento & Alto &
Estimadores de transformación ajustados exclusivamente sobre el
subconjunto de entrenamiento\\[2pt]
Confusión entre deriva temporal y cambio de versión & Medio &
Declaración explícita de la limitación; reporte por versión cuando el
volumen lo permita\\[2pt]
Memorización de configuraciones de mazo & Medio &
Protocolo C con separación por configuración\\[2pt]
Sesgo del estimador en turnos tempranos & Medio &
Reporte del soporte por ventana de turno junto a las métricas\\[2pt]
Descalibración del modelo seleccionado & Medio &
Medición previa y recalibración condicional\\
\bottomrule
\end{tabularx}
\end{table}

El riesgo de fuga por preprocesamiento merece mención explícita. Figueiredo
y Mendes señalan que, además de la fuga derivada de la similitud entre
observaciones, existen fuentes asociadas a la optimización de
hiperparámetros, la normalización de datos y el aumento de datos
\cite{figueiredo_leakage_2024}.

\razon{Esta fuente de fuga debe figurar en el registro de riesgos porque no
se detecta mediante inspección del código: un estimador de normalización
ajustado sobre la tabla completa antes de particionar produce un pipeline
funcional cuyos resultados están contaminados.}

\noindent\textbf{Expectativa sobre el efecto del protocolo.} Se anticipa
que la partición agrupada produzca métricas inferiores y mayor dispersión
entre ejecuciones que una partición aleatoria. Figueiredo y Mendes
documentan este comportamiento: las estrategias de menor fuga presentan
mayores pérdidas de entrenamiento y mayor varianza entre corridas
\cite{figueiredo_leakage_2024}. Dicho comportamiento corresponde al
funcionamiento correcto del protocolo y no a un defecto de su
calibración.

\razon{Esta declaración debe redactarse antes de disponer de resultados.
Ante métricas inferiores a las esperadas, la expectativa documentada
previamente protege la decisión de mantener el protocolo.}